\DocumentMetadata{
  pdfversion=2.0,
  pdfstandard=ua-2,
  lang=en-US,
  tagging=on,
  tagging-setup={math/setup=mathml-SE,tabsorder=structure}
}
\documentclass[11pt,letterpaper]{article}
\usepackage[margin=0.68in,includefoot]{geometry}
\usepackage{newcomputermodern}
\usepackage{amsmath,amscd,mathtools}
\usepackage{tikz,adjustbox}
\providecommand{\square}{\mdlgwhtsquare}
\usepackage[hidelinks]{hyperref}
\hypersetup{
  pdftitle={Combinatorics PhD Exam, May 2010},
  pdfauthor={University of Florida Department of Mathematics},
  pdfsubject={Graduate mathematics examination},
  pdfdisplaydoctitle=true,
  bookmarksopen=true
}
\setcounter{secnumdepth}{0}
\setlength{\parindent}{0pt}
\setlength{\parskip}{0.28em}
\setlength{\emergencystretch}{3em}
\setlength{\leftmargini}{2.15em}
\setlength{\leftmarginii}{2.65em}
\setlength{\leftmarginiii}{3em}
\renewcommand{\labelenumi}{\textbf{\theenumi.}}
\pagestyle{empty}
\raggedbottom
\allowdisplaybreaks
\newenvironment{examproblems}[1]{%
  \begin{enumerate}%
  \setcounter{enumi}{#1}%
  \setlength{\topsep}{0.35em}%
  \setlength{\partopsep}{0pt}%
  \setlength{\itemsep}{0.48em}%
  \setlength{\parsep}{0.18em}%
}{\end{enumerate}}
\newenvironment{examsubparts}{%
  \begin{enumerate}%
  \setlength{\topsep}{0.18em}%
  \setlength{\partopsep}{0pt}%
  \setlength{\itemsep}{0.16em}%
  \setlength{\parsep}{0.1em}%
}{\end{enumerate}}
\newenvironment{examinstructions}{%
  \par\begingroup\itshape
}{%
  \par\endgroup\vspace{0.18em}
}
\makeatletter
\renewcommand\section{\@startsection{section}{1}{0pt}%
  {-1.25ex plus -.25ex minus -.1ex}%
  {0.55ex plus .1ex}%
  {\normalfont\large\bfseries}}
\renewcommand{\@maketitle}{%
  \newpage\null\vskip -1em%
  {\footnotesize\bfseries
    \href{https://www.ufl.edu/}{University of Florida}, \href{https://math.ufl.edu/}{Department of Mathematics}\par}%
  \vskip -0.5em%
  \tagstructbegin{tag=Title}%
    {\LARGE\bfseries\href{https://gma.math.ufl.edu/exams/combinatorics-phd/}{Combinatorics PhD Exam}, May 2010\par}%
  \tagstructend%
  \vskip 0.25em%
  \tagmcbegin{artifact}%
    \hrule height 1pt%
  \tagmcend%
  \vskip 1em%
}
\makeatother
\title{Combinatorics PhD Exam, May 2010}
\author{}
\date{}
\begin{document}
\maketitle
\thispagestyle{empty}
\begin{examproblems}{0}
\item
This problem has two parts.
\begin{examsubparts}
\item[\textnormal{(a)}]
Prove that 
\[
\sum_{i=0}^{n} \binom{i}{k}=\binom{n+1}{k+1}.
\]
\item[\textnormal{(b)}]
Prove that the number of solutions of \(x_1+\cdots+x_k\leq n\) in positive integers is \(\binom{n}{k}\).
\end{examsubparts}
\item
This problem has two parts.
\begin{examsubparts}
\item[\textnormal{(a)}]
Prove that the number of Hamiltonian cycles in the complete bipartite graph \(K_{n,n}\) is \(\frac{1}{2}(n-1)!n!\).
\item[\textnormal{(b)}]
Let~\(e\) be an edge of the complete graph \(K_n\). Use Cayley’s Theorem to find the number of spanning trees in \(K_n-e\).
\end{examsubparts}
\item
Recall that the average degree of a graph~\(G\) of order~\(n\) is 
\[
av(G)=\frac{\sum_{v\in V}\deg(v)}{n}.
\]
\begin{examsubparts}
\item[\textnormal{(a)}]
Prove that 
\[
\chi(G)\leq 1+\Delta(G),
\]
 where~\(\chi\) denotes the chromatic number and~\(\Delta\) the maximum degree.
\item[\textnormal{(b)}]
For an arbitrary graph~\(G\), prove or give a counterexample to the statement 
\[
\chi(G)\leq 1+av(G).
\]
\end{examsubparts}
\item
Let~\(C\) be a binary code and let \(\overline{C}\) be the code obtained from~\(C\) by adding an even parity check. Find the minimum distance \(d(\overline{C})\) in terms of \(d(C)\).
\item
Prove that balanced incomplete block designs with the following parameters do not exist: 
\[
(12,8,6,4,3),\qquad (22,22,7,7,2).
\]
\item
Find the unique sequence \(\{a_n\}\) with 
\[
\sum_{k=0}^{n} a_k a_{n-k}=1.
\]
\item
Use inclusion-exclusion to find a formula for the number of surjective functions \(f:X\to Y\) with \(|X|=n\) and \(|Y|=m\).
\item
Let \(f(n,k)\) be the number of~\(k\)-subsets of \(\{1,2,\ldots,n\}\) that do not contain a pair of consecutive integers.
\begin{examsubparts}
\item[\textnormal{(a)}]
Find a recurrence for \(f(n,k)\) by considering the~\(k\)-subsets that do, and do not, contain~\(1\).
\item[\textnormal{(b)}]
Use part (a) to show that \(f(n,k)=\binom{n-k+1}{k}\).
\end{examsubparts}
\end{examproblems}
\end{document}
